\documentclass[11pt,a4paper]{article}

\usepackage[ngerman]{babel}
\usepackage[utf8]{inputenc}
\usepackage[T1]{fontenc}
\usepackage{geometry}
\usepackage{microtype}
\microtypesetup{expansion=false}
\emergencystretch=3em
\usepackage{tabularx}
\usepackage{array}
\usepackage{booktabs}
\usepackage{titlesec}
\usepackage{enumitem}
\usepackage{xcolor}
\usepackage{hyperref}

\geometry{a4paper, margin=2.5cm}

\definecolor{accent}{RGB}{40,60,90}

\titleformat{\section}{\normalfont\Large\bfseries\color{accent}}{\thesection}{1em}{}
\titleformat{\subsection}{\normalfont\large\bfseries\color{accent}}{\thesubsection}{1em}{}

\hypersetup{
    colorlinks=true,
    linkcolor=accent,
    urlcolor=accent,
    pdftitle={VR Climbing Speedrun - Game Design Document}
}

\setlist[itemize]{leftmargin=1.5em, itemsep=0.2em}

\begin{document}

% ---------------- Titelseite ----------------
\begin{titlepage}
    \centering
    \vspace*{2cm}

    {\Huge\bfseries\color{accent} VR Climbing Speedrun}\\[0.5em]
    {\Large Game Design Document}\\[2cm]

    {\large Virtual Reality: Konzeption und Anwendung immersiver Technologien}\\[0.3em]

    \vfill
    {\large \today}
\end{titlepage}

\tableofcontents
\newpage

% ---------------- 1. Projektübersicht ----------------
\section{Projektübersicht}

\subsection{Elevator Pitch}

Das Projekt ist ein VR-basiertes Climbing-Spiel, bei dem der Spieler eine vertikale Strecke so schnell wie möglich durchquert. Der Fokus liegt auf physikbasiertem Klettern und einem zeitbasierten Gameplay-Loop.

\subsection{Kernfantasie}

Der Spieler erlebt sich als VR-Kletterer, der durch präzise Handbewegungen eine Strecke effizient bewältigt und seine Zeit optimiert.

\subsection{Zielplattform}

\begin{itemize}
    \item VR-Headsets (Roomscale / 3DoF)
    \item Unity Engine
    \item XR Interaction Toolkit
\end{itemize}

% ---------------- 2. High Concept ----------------
\section{High Concept}

\subsection{Kernidee}

Der Spieler startet einen Run, klettert durch eine Low-Poly-Umgebung und versucht, das Ziel in möglichst kurzer Zeit zu erreichen. Das Gameplay basiert auf physikalischem Greifen und Gravitation.

\subsection{USPs}

\begin{itemize}
    \item Physikbasiertes VR-Klettern
    \item Kurze, wiederholbare Runs
    \item Skill-basiertes Movement durch Handinteraktion
    \item Direkte VR-Immersion ohne klassische Locomotion
\end{itemize}

% ---------------- 3. Core Gameplay ----------------
\section{Core Gameplay}

\subsection{Core Loop}

\begin{enumerate}
    \item Spieler startet im Main Menu
    \item Countdown (Ready -- Set -- Go)
    \item Timer startet
    \item Spieler klettert durch das Level
    \item Ziel wird erreicht
    \item Zeit wird angezeigt
    \item Run wird neu gestartet
\end{enumerate}

\subsection{Success State}

Ziel erreicht. Zeit wird gestoppt und angezeigt.

\subsection{Match Flow}

Countdown Phase $\rightarrow$ Active Run $\rightarrow$ Finish State $\rightarrow$ Restart Loop

% ---------------- 4. Game Rules ----------------
\section{Game Rules}

\subsection{Grundregeln}

\begin{itemize}
    \item Movement nur durch Klettern
    \item Gravitation ist aktiv
    \item Ziel ist die schnellste Zeit
    \item Keine klassischen Movement-Systeme vorhanden
\end{itemize}

\subsection{Interaktionen}

\begin{itemize}
    \item Greifen von Oberflächen
    \item Physikalisches Ziehen des Körpers
    \item Kletterbare Levelgeometrie
\end{itemize}

% ---------------- 5. Player System ----------------
\section{Player System}

\subsection{Fähigkeiten}

\begin{itemize}
    \item Greifen und Klettern
    \item Physikalisches Movement durch Handpositionen
    \item Kontrolle über vertikale Bewegung
\end{itemize}

\subsection{Input}

\begin{itemize}
    \item Grip Button: Klettern / Festhalten
    \item VR Controller: Handtracking
\end{itemize}

\subsection{Spielerwerte}

\begin{itemize}
    \item Kein Stamina-System implementiert
    \item Gravitation aktiv
    \item Movement abhängig von Handbewegung
\end{itemize}

\subsection{Feedback}

\begin{itemize}
    \item Timer UI
    \item Countdown UI
    \item Finish-Anzeige
    \item Audio Feedback über Musikstart beim Run
\end{itemize}

% ---------------- 6. VR Movement System ----------------
\section{VR Movement System}

\subsection{Climbing System}

\begin{itemize}
    \item Hände greifen Oberflächen
    \item Körper bewegt sich entgegengesetzt zur Handbewegung
    \item Bewegung basiert auf Hand-Delta
\end{itemize}

\subsection{Gravity System}

\begin{itemize}
    \item Permanente Gravitation aktiv
    \item Spieler fällt ohne Grip
    \item Kombination aus Klettern und Fallen möglich
\end{itemize}

\subsection{Wall Interaction}

\begin{itemize}
    \item Kanten können gelegentlich zu Hängenbleiben führen
    \item Kein aktives Push-Off System implementiert
    \item Level ist für Klettern optimiert
\end{itemize}

% ---------------- 7. Level Design ----------------
\section{Level Design}

\subsection{Struktur}

\begin{itemize}
    \item Vertikales Leveldesign
    \item Fokus auf Aufstieg
    \item Ein klarer Zielpunkt
\end{itemize}

\subsection{Environment}

\begin{itemize}
    \item Felsen
    \item Plattformen
    \item natürliche Kletterstrukturen
\end{itemize}

% ---------------- 8. UI / UX ----------------
\section{UI / UX}

\subsection{HUD}

\begin{itemize}
    \item Timer (laufend)
    \item Countdown (Ready / Set / Go)
    \item Finish Text mit Zeit
\end{itemize}

% ---------------- 9. Audio Design ----------------
\section{Audio Design}

Ein Musiktrack startet beim Run-Start. Musik endet/stoppt beim Finish.

% ---------------- 10. Technik ----------------
\section{Technik}

\subsection{Kernsysteme}

\begin{itemize}
    \item VR Climbing System
    \item Gravity System
    \item Timer System
    \item Scene Management
    \item Finish Detection
    \item UI System
\end{itemize}

\subsection{Architektur}

\begin{itemize}
    \item GameManager steuert Spielzustand
    \item FinishDetection triggert Run-Ende
    \item UI ist an XR Camera gebunden
    \item Scoreboard im Main Menu zeigt Bestzeiten
\end{itemize}

% ---------------- 11. Game Flow ----------------
\section{Game Flow}

\begin{center}
    \large
    Start $\rightarrow$ Countdown $\rightarrow$ Run $\rightarrow$ Finish $\rightarrow$ Time Display $\rightarrow$ Restart
\end{center}

% ---------------- 12. Evaluation ----------------
\section{Evaluation}

Das Spiel wurde getestet mit Fokus auf:

\begin{itemize}
    \item Usability (Steuerbarkeit)
    \item Immersion (Presence)
    \item Motion Comfort
    \item Gameplay Flow
\end{itemize}

% ---------------- 13. Final Statement ----------------
\section{Final Statement}

Das Projekt ist ein funktionaler VR Climbing Speedrun Prototyp. Alle Kernmechaniken (Movement, Gravity, Timer, Finish Loop und Scoreboard) sind implementiert und ermöglichen ein wiederholbares, zeitbasiertes Spielerlebnis.

\end{document}
